\chapter{Régression}\label{ch:iv1:regression}

\begin{ivtrains}
La régression telle qu'un chercheur quantitatif l'utilise~: ce que représentent réellement un coefficient et son erreur type, et ce qui se passe
sur les données financières, où les régresseurs sont bruités, où les observations se chevauchent, où la volatilité varie et où le signal est très faible. La plupart des questions demandent
un nombre que l'on puisse défendre, et les dernières reposent sur un effet que l'on pourrait facilement manquer dans son propre travail.
\end{ivtrains}

\begin{question}[$\star$ \iqroles{QR, QT, DS, Risque, Analyste} \iqfirm{toutes} \iqlevel{junior}]\label{iq:iv1:regression:1}
Vous régressez $y$ sur $x$ et obtenez une pente de 2. Quelles valeurs peut prendre la pente de la régression inverse, $x$ sur $y$~? Si la pente inverse vaut $0.18$, quelle est la corrélation de $x$ et de $y$~?
\end{question}

\begin{question}[$\star$ \iqroles{QR, QT, DS, Risque} \iqfirm{toutes} \iqlevel{junior}]\label{iq:iv1:regression:2}
Vous recopiez chaque ligne d'un jeu de données, ce qui le rend deux fois plus long, puis vous refaites la même régression linéaire. Que deviennent les coefficients, leurs erreurs types, les statistiques $t$ et $R^2$~? Pourquoi cela compte-t-il lorsque les observations se chevauchent dans le temps~?
\end{question}

\begin{question}[$\star$ \iqroles{QR, QT, Trader, Risque} \iqfirm{toutes} \iqlevel{junior}]\label{iq:iv1:regression:3}
Une action a une volatilité de 30~\% par an et un indice une volatilité de 15~\%, avec une corrélation de $0.8$. Combien de dollars de l'indice faut-il vendre à découvert par dollar d'action pour que la variance de la position couverte soit la plus petite possible~? Quelle volatilité reste-t-il~? Que se passe-t-il si l'on prend le rapport issu de la régression de l'indice sur l'action, puis qu'on l'inverse~?
\end{question}

\begin{question}[$\star$ \iqroles{QR, DS, Risque, Analyste} \iqfirm{toutes} \iqlevel{junior}]\label{iq:iv1:regression:4}
Une régression avec une constante donne un $R^2$. Quelles valeurs peut-il prendre, et que signifie $R^2 = 1$~? Que se passe-t-il lorsqu'on ajoute un régresseur supplémentaire~? Si l'on supprime la constante mais que l'on calcule quand même $1 - \text{RSS}/\sum (y_i - \bar y)^2$, où RSS est la somme des carrés des résidus, que peut-il arriver~?
\end{question}

\begin{question}[$\star$ \iqroles{QR, QT, Analyste, Risque} \iqfirm{hedge fund, banque} \iqlevel{junior}]\label{iq:iv1:regression:5}
Vous régressez les rendements d'un fonds sur ceux du marché et estimez un alpha, l'ordonnée à l'origine de la régression, de 2~\% par an. Les résidus ont une volatilité de 10~\% par an, et vous disposez de 4 ans de données. Quelle est la statistique $t$ de l'alpha~? Si l'alpha vrai vaut 2~\%, combien d'années de données faudrait-il pour obtenir une statistique $t$ égale à 2~?
\end{question}

\begin{question}[$\star$ \iqroles{QR, QT, DS} \iqfirm{hedge fund, teneur de marché, crypto} \iqlevel{junior}]\label{iq:iv1:regression:6}
Vous régressez les rendements quotidiens sur un signal standardisé, à partir de 2\,500 jours de données, et obtenez $R^2 = 0.5~\%$. Quelle est la statistique $t$ de la pente~? Ce signal a-t-il une valeur~? Estimez le ratio de Sharpe annuel d'une position proportionnelle au signal, le ratio de Sharpe étant le rendement moyen divisé par son écart type.
\end{question}

\begin{question}[$\star\star$ \iqroles{QR, QD, DS, Risque} \iqfirm{toutes} \iqlevel{junior}]\label{iq:iv1:regression:7}
Dérivez à la main l'estimateur des moindres carrés de $\beta$ dans $y = X\beta + \varepsilon$. Quand existe-t-il~? Pourquoi ne le calcule-t-on pas en inversant $X^\top X$~? Comment l'estimer lorsque $X$ est trop grand pour tenir en mémoire~?
\end{question}

\begin{question}[$\star\star$ \iqroles{QR, QT, DS, Risque} \iqfirm{toutes} \iqlevel{junior, confirmé}]\label{iq:iv1:regression:8}
Vous voulez le coefficient de $x$ dans une régression de $y$ sur $x$ et $z$. La procédure A régresse $y$ sur $z$, puis régresse les résidus ainsi obtenus sur $x$. La procédure B régresse $y$ sur $z$ et $x$ sur $z$, puis régresse le premier jeu de résidus sur le second. Laquelle donne le même coefficient que la régression complète~? Si $x$ et $z$ ont une corrélation de $0.8$ et que le coefficient vrai de $x$ vaut 1, que donne l'autre procédure~?
\end{question}

\begin{question}[$\star\star$ \iqroles{QR, QT, DS, Risque} \iqfirm{toutes} \iqlevel{junior}]\label{iq:iv1:regression:9}
Le modèle vrai est $y = x - 3z + \varepsilon$, où $x$ et $z$ ont une variance unitaire et une corrélation de $0.5$, et où $\varepsilon$ est un bruit indépendant. Vous n'observez pas $z$ et régressez $y$ sur $x$ seul. Quelle pente obtenez-vous asymptotiquement~?
\end{question}

\begin{question}[$\star\star$ \iqroles{QR, QT, DS, Risque} \iqfirm{toutes} \iqlevel{junior, confirmé}]\label{iq:iv1:regression:10}
La relation vraie est $y = x^* + \varepsilon$. Vous observez $x = x^* + n$, où $x^*$ et le bruit de mesure $n$ ont chacun une variance de 1 et sont indépendants. Asymptotiquement, quelle pente obtient-on en régressant $y$ sur $x$~? Et si le bruit de mesure porte sur $y$ et non sur $x$~? Si les deux variables sont bruitées, comment encadrer la vraie pente à l'aide de deux régressions~?
\end{question}

\begin{question}[$\star\star$ \iqroles{QR, DS, Risque} \iqfirm{toutes} \iqlevel{junior}]\label{iq:iv1:regression:11}
Deux régresseurs ont une corrélation de $0.9$. De quel facteur l'erreur type de chaque coefficient est-elle plus grande que si les deux variables étaient non corrélées~? Qu'est-ce qui reste bien estimable~? Peut-il arriver que les deux statistiques $t$ soient inférieures à 2 alors que le test $F$ conjoint est très significatif~?
\end{question}

\begin{question}[$\star\star$ \iqroles{QR, QT, DS} \iqfirm{toutes} \iqlevel{junior}]\label{iq:iv1:regression:12}
Un point est tiré uniformément dans le triangle de sommets $(0,0)$, $(1,0)$ et $(1,1)$. Déterminez la droite de régression de population de $y$ sur $x$, ainsi que la droite de régression de $x$ sur $y$. Quelle est la corrélation de $x$ et de $y$~?
\end{question}

\begin{question}[$\star\star$ \iqroles{QR, DS, QD} \iqfirm{toutes} \iqlevel{junior, confirmé}]\label{iq:iv1:regression:13}
Vous incluez par mégarde deux fois la même colonne dans une régression. Que fait la méthode des moindres carrés ordinaires~? Que fait la régression ridge~? Que fait le lasso~?
\end{question}

\begin{question}[$\star\star$ \iqroles{QR, DS} \iqfirm{toutes} \iqlevel{junior, confirmé}]\label{iq:iv1:regression:14}
Les colonnes de $X$ sont orthonormées, de sorte que $X^\top X = I$, et les coefficients des moindres carrés ordinaires valent $3$, $-0.5$ et $1.2$. Déterminez les coefficients qui minimisent $\tfrac12\lVert y - Xb\rVert^2 + \lambda\lVert b\rVert_1$ (lasso), puis ceux qui minimisent $\tfrac12\lVert y - Xb\rVert^2 + \tfrac{\lambda}{2}\lVert b\rVert_2^2$ (régression ridge), avec $\lambda = 1$. Dérivez la règle que vous utilisez.
\end{question}

\begin{question}[$\star\star$ \iqroles{QR, QT, DS, Risque, Analyste} \iqfirm{toutes} \iqlevel{junior}]\label{iq:iv1:regression:15}
Sous quelles hypothèses les coefficients d'une régression linéaire sont-ils sans biais, et sous quelles hypothèses supplémentaires leurs erreurs types sont-elles correctes~? Sur les rendements quotidiens d'actions, lesquelles de ces hypothèses échouent, et que faites-vous dans chaque cas~?
\end{question}

\begin{question}[$\star\star$ \iqroles{QR, QT, DS} \iqfirm{toutes} \iqlevel{junior}]\label{iq:iv1:regression:16}
Vous régressez $y$ sur 20 régresseurs avec 100 observations et une constante, alors qu'en réalité $y$ est indépendant de tous. Quel est le $R^2$ espéré~? Quel est le $R^2$ ajusté espéré~? Combien de ces 20 coefficients devriez-vous attendre avec $|t| > 2$~?
\end{question}

\begin{question}[$\star\star$ \iqroles{QR, QT, Risque, DS} \iqfirm{banque, hedge fund} \iqlevel{junior}]\label{iq:iv1:regression:17}
Vous générez deux marches aléatoires indépendantes, longues chacune de 250 pas, et régressez l'une sur l'autre en utilisant les niveaux. À quelle fréquence la pente est-elle «~significative~» avec $|t| > 1.96$, et quel $R^2$ observe-t-on en moyenne~? Qu'est-ce qui change avec 1\,000 pas~? Que faut-il faire avec des séries de prix à la place~?
\end{question}

\begin{question}[$\star\star$ \iqroles{QR, QT, DS} \iqfirm{hedge fund, teneur de marché, banque} \iqlevel{junior, confirmé}]\label{iq:iv1:regression:18}
Vous disposez de 500 signaux de rendement faibles et corrélés, sur 5 ans de données quotidiennes, soit environ 1\,250 observations. Ajusteriez-vous une régression ridge ou un lasso~? Qu'est-ce qui vous ferait changer d'avis~?
\end{question}

\begin{question}[$\star\star$ \iqroles{QR, QT, Risque, Trader} \iqfirm{banque, hedge fund} \iqlevel{junior, confirmé}]\label{iq:iv1:regression:19}
Le rendement quotidien vrai d'une action vaut exactement $1.2$ fois celui du marché. Mais elle ne se traite que sur une moitié des jours tirée au hasard, et les autres jours son cours reste simplement égal au dernier cours traité. Vous régressez son rendement quotidien observé sur le rendement du marché du même jour. Quelle pente obtenez-vous~? Comment retrouver $1.2$~?
\end{question}

\begin{question}[$\star\star\star$ \iqroles{QR, DS, Risque} \iqfirm{toutes} \iqlevel{junior, confirmé}]\label{iq:iv1:regression:20}
Soit $x$ une variable normale centrée réduite et $y = \beta x + |x|\,\varepsilon$, avec $\varepsilon$ une variable normale centrée réduite indépendante de $x$, de sorte que le bruit est grand lorsque $|x|$ est grand. Pour un grand échantillon de taille $n$, quelle est la vraie erreur type de la pente MCO, et que donne la formule usuelle~? Comment corriger les erreurs types, et comment utiliser cela pour améliorer l'estimation~?
\end{question}

\begin{question}[$\star\star\star$ \iqroles{QR, QT, DS} \iqfirm{hedge fund, teneur de marché} \iqlevel{junior, confirmé}]\label{iq:iv1:regression:21}
Les rendements quotidiens sont indépendants, et un signal n'a aucun lien avec eux. Vous régressez le rendement total des 20 jours suivants sur le signal du jour, en échantillonnant chaque jour. De quel facteur les erreurs types usuelles sous-estiment-elles la vérité si le signal est (a) indépendant d'un jour à l'autre, (b) presque constant sur 20 jours, (c) un processus autorégressif d'ordre 1 de coefficient $0.9$~?
\end{question}

\begin{question}[$\star\star\star$ \iqroles{QR, DS, QT} \iqfirm{toutes} \iqlevel{junior, confirmé}]\label{iq:iv1:regression:22}
$y$ a une corrélation de $0.6$ avec $x_1$ et une corrélation de $0.6$ avec $x_2$, mais vous ne connaissez pas la corrélation entre $x_1$ et $x_2$. Quelles valeurs peut prendre le $R^2$ de la régression de $y$ sur les deux~? Considérons ensuite un autre cas~: $x_1$ n'est pas corrélé à $y$, $x_2$ a une corrélation de $0.4$ avec $y$, et $x_1$ et $x_2$ ont une corrélation de $0.9$. $x_1$ est-il inutile~?
\end{question}

\begin{question}[$\star\star\star$ \iqroles{QR, QT} \iqfirm{hedge fund, teneur de marché, crypto} \iqlevel{confirmé}]\label{iq:iv1:regression:23}
Deux signaux ont chacun une corrélation de $0.05$ avec le rendement du lendemain. Quelle corrélation la meilleure combinaison linéaire des deux peut-elle atteindre si leur corrélation mutuelle vaut $0.5$, $0$ ou $-0.5$~? Dans le cas de $0.5$, supposez que la corrélation du second signal avec le rendement soit de $0.02$ et non de $0.05$. Quel signe doit avoir son poids~?
\end{question}

\begin{question}[$\star\star\star$ \iqroles{QR, DS} \iqfirm{toutes} \iqlevel{confirmé}]\label{iq:iv1:regression:24}
Une régression comporte un seul régresseur avec $x^\top x = 1$, et un bruit de variance $\sigma^2 = 1$. Le coefficient vrai vaut $\beta = 0.5$. Les MCO donnent $\hat\beta$, et la régression ridge donne $\hat\beta/(1 + \lambda)$. Quel $\lambda$ minimise l'erreur quadratique moyenne de l'estimateur, et de combien cette erreur est-elle inférieure à celle des MCO~? Qu'est-ce que cela dit du rétrécissement pour les signaux faibles~?
\end{question}

\begin{question}[$\star\star\star$ \iqroles{QR, QT, Risque, DS} \iqfirm{toutes} \iqlevel{junior, confirmé}]\label{iq:iv1:regression:25}
$x$ et $y$ sont conjointement normales, de variance unitaire et de corrélation $0.5$. Vous ne gardez que les observations avec $x > 0$ et régressez $y$ sur $x$. Quelles pente et corrélation obtenez-vous~? Quelle pente obtenez-vous si vous ne gardez à la place que les observations avec $y > 0$~?
\end{question}

\begin{question}[$\star\star\star$ \iqroles{QR, DS, Risque} \iqfirm{banque, hedge fund} \iqlevel{junior, confirmé}]\label{iq:iv1:regression:26}
Une régression de $y$ sur $x$ avec constante comporte 1\,000 observations, où $x$ est de moyenne 0 et $\sum x_i^2 = 1000$. Une observation a $x = 10$ et un résidu de 2. Sans réajuster le modèle, déterminez son levier, la variation de la pente si on la retire, et l'erreur de prévision qui lui correspond dans un ajustement sans elle.
\end{question}

\begin{question}[$\star\star\star$ \iqroles{QR, QT, DS} \iqfirm{hedge fund, banque, teneur de marché} \iqlevel{confirmé}]\label{iq:iv1:regression:27}
Le modèle vrai est $y = x + u$, et $x = z + u + v$, où $z$, $u$ et $v$ sont indépendantes et de variance unitaire. Ainsi $x$ est corrélé à l'erreur $u$. Quelle pente donnent asymptotiquement les MCO de $y$ sur $x$~? Comment utiliser $z$ pour retrouver la vraie pente~? Qu'est-ce qui ne va pas si $z$ n'est que faiblement lié à $x$~?
\end{question}

\begin{question}[$\star\star\star$ \iqroles{QR, QT, DS} \iqfirm{hedge fund, teneur de marché} \iqlevel{junior, confirmé}]\label{iq:iv1:regression:28}
Vous ajustez une régression MCO à 20 régresseurs sur 1\,000 observations, et le $R^2$ vrai du meilleur prédicteur linéaire vaut 1~\%. Quel $R^2$ hors échantillon doit-on attendre approximativement, mesuré par rapport à la prévision de la vraie moyenne~? De combien de données aurait-on besoin pour que le $R^2$ hors échantillon atteigne zéro~?
\end{question}

\begin{question}[$\star\star\star\star$ \iqroles{QR, QT} \iqfirm{hedge fund} \iqlevel{confirmé}]\label{iq:iv1:regression:29}
Considérons $r_{t+1} = \beta x_t + u_{t+1}$ et $x_{t+1} = \rho x_t + v_{t+1}$, où $\beta = 0$, $\rho = 0.95$, les chocs $u$ et $v$ ont une variance unitaire et une corrélation de $-0.9$, et vous disposez de 120 observations. Vous régressez $r_{t+1}$ sur $x_t$ avec une constante. Quelle est approximativement la moyenne de la pente estimée~? Vous testez $\beta = 0$ contre $\beta > 0$ avec la statistique $t$ usuelle au seuil nominal de 5~\%. Le test rejette-t-il plus ou moins souvent que 5~\%, et à peu près combien de fois~?
\end{question}

\begin{question}[$\star\star\star\star$ \iqroles{QR, QT, DS} \iqfirm{hedge fund, banque} \iqlevel{confirmé}]\label{iq:iv1:regression:30}
$y = x_1 + \varepsilon$, donc seul $x_1$ compte. Vous observez $w = x_1 + n$, une version bruitée de $x_1$ avec une variance de bruit de $0.5$, ainsi qu'une seconde variable $x_2$. $x_1$ et $x_2$ ont toutes deux une variance unitaire, leur corrélation vaut $0.8$, et $x_2$ n'a aucun effet sur $y$ une fois $x_1$ connu. Vous régressez $y$ sur $w$ et $x_2$. Quels coefficients obtenez-vous asymptotiquement~? Que concluriez-vous sur $x_2$, et quelle est l'erreur~?
\end{question}
